\documentclass[12pt]{article}

% -----------------------------------------------
% PACOTES OBRIGATÓRIOS (SBC TEMPLATE OFICIAL)
% -----------------------------------------------
\usepackage{sbc-template}           % Estilo oficial SBC — inclui geometry, times, caption e sections
\usepackage{graphicx,url}
\usepackage[T1]{fontenc}
\usepackage[utf8]{inputenc}
\usepackage[brazil]{babel}

% Pacotes complementares (não conflitam com sbc-template)
\usepackage{enumitem}
\setlist[itemize]{label=--}
\usepackage{booktabs}
\usepackage{tikz}
\usetikzlibrary{shapes.geometric,arrows.meta,positioning,fit,backgrounds}

\sloppy

% -----------------------------------------------
% CABEÇALHO
% -----------------------------------------------
\title{Medição Automatizada de QoE e UX com Playwright:\\
  Um Instrumento de Auditoria Passiva em Aplicações Web}

\author{Mariana Ribeiro\inst{1}, Vinícius Schineider Januário Viana\inst{1},\\
        Simone Azevedo Bandeira de Melo Aquino\inst{1}}

\address{Instituto Federal de Educação, Ciência e Tecnologia do Maranhão (IFMA) -- Campus Imperatriz\\
  Imperatriz -- MA -- Brasil
  \email{\{mariana.r,viniciusschineider\}@acad.ifma.edu.br,
  simonebandeira@ifma.edu.br}
}

\begin{document}

\maketitle

% -----------------------------------------------
% RESUMOS — sbc-template.sty insere "Abstract."
% e "Resumo." automaticamente em itálico negrito;
% NÃO repetir o prefixo aqui.
% -----------------------------------------------
\begin{abstract}
This paper proposes and evaluates (proof of concept) a Playwright-based passive
auditing instrument for QoE/UX in web applications. A \textit{listen-only} script
captures \textit{dead clicks}, long hesitations ($\geq 10$ s), and flow duration.
In a controlled mixed-methods study ($N = 9$) on DevShare (four flows), logs
were triangulated with UEQ-S and open-ended responses. 78\% reported missing
visual feedback, consistent with recorded \textit{dead click} events. Power
analysis indicates $N \geq 30$ for confirmatory claims.
\end{abstract}

\begin{resumo}
Este artigo propõe e avalia (como prova de conceito) um instrumento baseado em
Playwright para auditoria passiva de QoE/UX. Um script \textit{listen-only}
captura \textit{dead clicks}, hesitações longas ($\geq 10$ s) e duração de
fluxos, sem instrumentação intrusiva no código de produção. Em um estudo
controlado de métodos mistos ($N = 9$) na plataforma DevShare (quatro fluxos),
triangulamos logs com UEQ-S e respostas abertas. 78\% relataram ausência de
\textit{feedback} visual, consistente com os \textit{dead clicks} registrados.
A análise de poder sugere $N \geq 30$ para conclusões confirmatórias.
\end{resumo}

% ==========================================
% CORPO DO ARTIGO
% ==========================================
\section{Introdução}

A avaliação da Qualidade de Experiência (QoE) e da Usabilidade (UX) em aplicações Web tem sido, historicamente, conduzida por meio de métodos manuais e subjetivos. Os métodos atuais dependem majoritariamente de intervenções aplicadas a posteriori ou de práticas de \textit{Quality Assurance} (QA) focadas na validação sintética de regressões funcionais \cite{leppanen2024}. No entanto, a lacuna metodológica na automação de testes reside na distinção clássica entre sintaxe e semântica: algoritmos podem validar códigos sem possuir intencionalidade ou compreensão sobre a frustração emocional vivenciada durante a interação \cite{searle1982,rousi2019}.

Além do aspecto humano, a ineficiência no design de interfaces pode induzir comportamentos erráticos --- como cliques repetidos (\textit{dead clicks}) e hesitações prolongadas. Em aplicações Web, tais padrões podem estar associados a retrabalho (ações repetidas e fluxos reiniciados) e, potencialmente, a processamento e tráfego redundantes \cite{oyedeji2019,rani2024}.

Para endereçar este problema, este artigo propõe um instrumento (um arcabouço técnico de auditoria passiva) para medir QoE/UX de forma automatizada durante a interação. O trabalho altera a aplicação convencional do \textit{Playwright}: em vez de executar rotinas sintéticas de clique, a ferramenta atua como observador passivo da jornada de usuários reais, capturando rastros de telemetria sem instrumentação intrusiva.

O presente artigo detalha a arquitetura desta proposta e apresenta uma prova de conceito (PoC) em estudo controlado na plataforma Web ``DevShare''. O estudo investiga, de modo exploratório, como sinais sintáticos capturados pela máquina (tempos e eventos) \cite{rousi2019} se alinham à percepção subjetiva reportada via UEQ-S e a evidências qualitativas.

Em síntese, este trabalho contribui com (i) um instrumento de auditoria passiva, baseado em Playwright, para captura não intrusiva de sinais comportamentais associados a fricção (\textit{dead clicks}, hesitações longas e duração de fluxos); (ii) um protocolo de estudo controlado de métodos mistos, com coleta objetiva e subjetiva, para triangulação dos sinais capturados; e (iii) uma prova de conceito em uma aplicação Web realista (DevShare), evidenciando alinhamento inicial entre registros automatizados e relatos qualitativos.

Para tal, este artigo está organizado da seguinte forma: a Seção 2 apresenta a fundamentação teórica e trabalhos relacionados; a Seção 3 descreve a arquitetura do instrumento e o desenho metodológico do estudo; a Seção 4 reporta os resultados e discute os achados a partir da triangulação quali-quanti; por fim, a Seção 5 apresenta conclusões, limitações e trabalhos futuros.

\section{Fundamentação Teórica e Trabalhos Relacionados}

Esta seção estabelece os conceitos necessários para interpretar a proposta e os resultados. Em particular, delimita-se o papel da máquina como coletora de \textit{sintaxe} de interação e a necessidade de triangulação com dados humanos para inferência de \textit{semântica} (percepção).

O instrumento proposto fundamenta-se na intersecção entre a filosofia da percepção e a automação de medição de QoE/UX.

A experiência do usuário divide-se fundamentalmente em metas pragmáticas e hedônicas \cite{grigis2024}, cujas barreiras são sintetizadas na Tabela 1.

\begin{table}[htbp]
  \centering
  \caption{Diretores e barreiras por meta de UX}
  \begin{tabular}{p{2.5cm} p{4cm} p{7cm}}
    \toprule
    \textbf{Metas de UX} & \textbf{Diretores (Drivers)} & \textbf{Barreiras (Barriers)} \\
    \midrule
    \textbf{Pragmáticas} & Utilidade & Limitações técnicas, disponibilidade, perda econômica \\
    \textbf{Hedônicas} & Surpresa, fascinação & Ameaça à identidade, política (filtros e tabus) \\
    \bottomrule
  \end{tabular}
\end{table}

Ferramentas clássicas de QA atuam prioritariamente sobre barreiras pragmáticas. Conforme o experimento mental do ``Quarto Chinês'' \cite{searle1982}, um script de automação manipula a interface corretamente através de regras pré-definidas, mas carece de agência para interpretar a degradação da experiência. Sem o acompanhamento empírico humano, os logs automatizados tendem a atuar como ``significantes vazios'' \cite{rousi2019}. A abordagem aqui proposta delega à máquina (Playwright) exclusivamente a captura da \textit{sintaxe} das interações em tempo real (tempos e eventos), condicionando a interpretação \textit{semântica} ao cruzamento com dados humanos (UEQ-S).

No que diz respeito à automação passiva de QoE/UX, estudos sustentam a viabilidade de afastar a dependência exclusiva da subjetividade humana em avaliações. No cenário nacional, Garcia et al. (2023) já apontavam a necessidade de estruturas para a automatização do diagnóstico de UX, destacando a urgência de métricas quantitativas. Internacionalmente, Keshvadi e Williamson (2020) empregaram a ferramenta MoVIE para capturar tráfego passivamente e inferir QoE, enquanto Banaouas (2019) evidenciou que 55\% das heurísticas de usabilidade podem ser diagnosticadas via sequenciamento de eventos de GUI.

Como implicação secundária, estudos recentes documentam a interação entre usabilidade e aspectos de eficiência de recursos em aplicações Web \cite{oyedeji2019,rani2024}. A premissa adotada é que a eliminação de atritos de interface pode reduzir ações repetitivas, que se traduziriam em processamento e tráfego redundantes. Neste contexto, rastrear falhas de usabilidade via ferramentas de escuta passiva --- estruturalmente mais leves que automações baseadas em processamento de imagens, como o Sikuli \cite{savoldi2018} --- constitui uma tática de prevenção de ineficiência do produto.

\section{Metodologia e Arquitetura do Instrumento}

Nesta seção, descreve-se a proposta (instrumento) e o desenho do estudo utilizado para avaliá-la como prova de conceito. A pesquisa caracteriza-se como um estudo metodológico e uma Prova de Conceito (PoC), operacionalizada por meio de um desenho quase-experimental controlado. O objetivo é investigar a associação empírica entre falhas de interface (sintaxe), percepção do usuário (semântica) e sinais objetivos de fricção observáveis durante a execução de tarefas.

\subsection{O Playwright como Auditor Passivo no CI/CD}

A principal característica da arquitetura proposta é a instanciação do \textit{Playwright} em modo \textit{listen-only} (escuta). Idealizado para ser acoplado ao pipeline de homologação (CI/CD), a ferramenta injeta um script ouvinte via API \texttt{page.addInitScript}, interceptando eventos customizados \texttt{qoe:step} no \textit{front-end}. O sistema monitora três indicativos:
 (i) a duração do fluxo (tempo total de atividade para conclusão da tarefa); (ii) \textit{dead clicks} (cliques sem resposta perceptível do sistema, capturados via instrumentação de eventos DOM); e (iii) hesitações (períodos de inatividade superiores a 10 segundos entre transições de página, limiar definido a partir da literatura de tempos de resposta aceitáveis em interfaces Web \cite{nielsen1994}).

Esta arquitetura viabiliza uma coleta não intrusiva, dispensando a inserção de bibliotecas de telemetria de terceiros no código de produção.\footnote{Artefatos do estudo (código-fonte, dados anonimizados e scripts de análise): \texttt{[URL\_ANONIMA]}.}

\begin{figure}[htbp]
\centering
\begin{tikzpicture}[
  font=\small,
  box/.style={rectangle, draw, rounded corners=3pt, minimum width=2.2cm,
              minimum height=0.9cm, align=center, fill=white},
  grp/.style={rectangle, draw=gray!60, rounded corners=4pt,
              inner sep=6pt, fill=gray!8},
  arr/.style={-{Stealth[length=5pt]}, thick},
  lbl/.style={font=\tiny, midway, above, sloped}
]

% --- Nós ---
\node[box] (user) {Usuário\\(participante)};

% Grupo: DevShare
\node[box, right=1.1cm of user]  (iface)  {Interface\\React/Next.js};
\node[box, below=0.5cm of iface] (event)  {\texttt{qoe:step}\\dispatcher};
\begin{scope}[on background layer]
  \node[grp, fit=(iface)(event), label={[font=\footnotesize]above:DevShare · localhost}] (fe) {};
\end{scope}

% Grupo: Playwright
\node[box, right=1.1cm of iface]  (init)    {\texttt{addInitScript}\\(ouvinte)};
\node[box, below=0.5cm of init]   (capture) {Captura DOM\\+ timestamps};
\begin{scope}[on background layer]
  \node[grp, fit=(init)(capture), label={[font=\footnotesize]above:Playwright · listen-only}] (pw) {};
\end{scope}

% Grupo: Saída
\node[box, right=1.1cm of init]   (dc)   {\textit{dead clicks}};
\node[box, below=0.25cm of dc]    (hes)  {hesitações $\geq$10\,s};
\node[box, below=0.25cm of hes]   (dur)  {duração do fluxo};
\begin{scope}[on background layer]
  \node[grp, fit=(dc)(hes)(dur), label={[font=\footnotesize]above:Saída · CSV}] (out) {};
\end{scope}

% --- Setas ---
\draw[arr] (user)    -- node[lbl]{interage} (iface);
\draw[arr] (iface)   -- (event);
\draw[arr] (event.east) -- ++(0.4,0) |- node[lbl, near end]{escutado por} (init.west);
\draw[arr] (init)    -- (capture);
\draw[arr] (capture.east) -- ++(0.4,0) |- (dc.west);
\draw[arr] (capture.east) -- ++(0.4,0) |- (hes.west);
\draw[arr] (capture.east) -- ++(0.4,0) |- (dur.west);

\end{tikzpicture}
\caption{Arquitetura da auditoria passiva via Playwright (\texttt{addInitScript} e eventos \texttt{qoe:step}).}
\label{fig:arquitetura}
\end{figure}

\subsection{Objeto de Estudo: A Plataforma DevShare}

Os testes empregam a ``DevShare'', aplicação Web desenvolvida como artefato de pesquisa para o nicho de compartilhamento de dicas entre profissionais de tecnologia (desenvolvedores, designers e operações). Construída com Next.js (App Router) e React, sua interface é projetada para avaliações de tarefas pragmáticas e utilitárias.\footnote{Instância pública: \texttt{[URL\_VERCEL]} (SPA; persistência em \texttt{localStorage}).}

A fim de isolar instabilidades inerentes a conexões externas de internet, a aplicação utiliza persistência em memória local (\texttt{localStorage}). Para reproduzir latências de rede sem dependência de infraestrutura externa, o código implementa mocks assíncronos (via \texttt{setTimeout} e resolução de promessas), que introduzem atrasos programados nas operações de escrita e leitura. O Playwright captura os efeitos comportamentais dessas latências --- hesitações e \textit{dead clicks} --- via eventos DOM e timestamps de navegação, independentemente do mecanismo subjacente de persistência. Esta escolha de design isola a variável de interesse (atrito de \textit{feedback} da interface) dos fatores de variabilidade de infraestrutura.


\subsection{Desenho Experimental e Cenários}

A execução das tarefas compreende quatro fluxos independentes, mapeados na Tabela 2. Os fluxos foram definidos para cobrir as principais jornadas pragmáticas da plataforma: exploração de conteúdo, criação de publicação, interação social e configuração de perfil.

\begin{table}[htbp]
  \centering
  \caption{Fluxos e condições experimentais}
  \begin{tabular}{p{2.8cm} p{2.5cm} p{2cm} p{5.2cm}}
    \toprule
    \textbf{Fluxo} & \textbf{Tipo} & \textbf{Condição} & \textbf{Característica} \\
    \midrule
    \textbf{explorar-salvar}  & Exploratório & Baseline  & Navegação livre pelo feed; maior volume de conteúdo \\
    \textbf{postar-dica}      & Criação      & Degradado & Ausência de \textit{feedback} visual após submissão do formulário \\
    \textbf{criar-perfil}     & Configuração & Degradado & Ausência de indicador de carregamento após confirmação \\
    \textbf{comentar-reagir}  & Social       & Neutro     & Interação com publicações existentes; sem degradação induzida \\
    \bottomrule
  \end{tabular}
\end{table}

Nos fluxos degradados, a supressão de \textit{feedback} de carregamento configura uma violação das heurísticas de visibilidade do estado do sistema \cite{nielsen1994,molich1990}, com o intuito de induzir comportamentos erráticos mensuráveis --- em especial \textit{dead clicks} e hesitações prolongadas. Os participantes executaram as tarefas sem conhecimento prévio das degradações inseridas.

\subsection{Ambiente de Teste e Protocolo de Coleta de Dados}

O experimento contemplou 9 participantes com perfil de profissionais ou estudantes de tecnologia. O desenho é within-subjects: cada participante executou os quatro fluxos sequencialmente, em ordem fixa, em uma única sessão. A ordem foi padronizada para controlar efeitos de aprendizado: (1) explorar-salvar, (2) postar-dica, (3) criar-perfil, (4) comentar-reagir. A coleta foi presencial, em computador controlado de laboratório, para mitigar variabilidade de hardware, seguindo três etapas: (i) briefing, no qual o participante recebeu o roteiro de tarefas e assinou o Termo de Consentimento Livre e Esclarecido (TCLE), sem conhecimento das degradações experimentais; (ii) execução interativa, com a plataforma servida em \texttt{localhost} e registro em segundo plano, pelo Playwright, de timestamps de navegação, contagem de \textit{dead clicks} e duração de cada fluxo; e (iii) avaliação subjetiva ao fim das tarefas, via \textit{User Experience Questionnaire -- Short} (UEQ-S) e três perguntas abertas (Q1--Q3).\footnote{UEQ-S e planilha de pontuação: \texttt{ueq-online.org}.}

A dissociação temporal metodológica entre a coleta mecânica (silenciosa) e o relato humano permite a estruturação de um estudo de métodos mistos, viabilizando a triangulação dos dados.

\subsection{Procedimentos de Análise dos Dados e Métricas}\label{subsec:analise-metricas}

Os dados coletados foram analisados de forma descritiva e correlacional, combinando sinais objetivos (logs) e evidências subjetivas (UEQ-S e respostas abertas). Adotaram-se três níveis de análise: (i) por fluxo (agregado), com duração média do fluxo e contagem média de hesitações longas; (ii) por participante (sessão), com total de \textit{dead clicks}, total de hesitações longas e duração total da sessão; e (iii) por evidência qualitativa, com categorias temáticas associadas a falhas de \textit{feedback} e navegação.

Para consolidar os sinais objetivos em um indicador único de sessão, definiu-se um \textit{FrictionScore} composto em dois estágios. Primeiro, para cada participante, calculou-se um score de sessão ponderando \textit{dead clicks} (peso 0,4), hesitações longas $\geq 10$ s (peso 0,4) e duração total (peso 0,2), com normalização min-max por participante. Segundo, para comparação entre fluxos, agregou-se a média por fluxo e reaplicaram-se pesos de duração (0,6) e hesitações (0,4), resultando em um índice de 0 a 100. Os pesos refletem a premissa de que eventos discretos de clique e pausas longas são sinais mais diretos de fricção do que a duração bruta, que pode ser legitimamente alta em fluxos exploratórios.

Na análise quantitativa, reportaram-se estatísticas descritivas e aplicou-se a correlação de Spearman entre o \textit{FrictionScore} de sessão e as escalas pragmática e hedônica do UEQ-S. Adicionalmente, realizou-se análise de poder \textit{a priori} via transformação Fisher-Z com $|\rho| = 0{,}5$, $\alpha = 0{,}05$ e poder $= 0{,}80$.

Na análise qualitativa, as respostas abertas (Q1--Q3) foram submetidas a análise de conteúdo com categorização dedutiva \cite{bardin1977}. As categorias foram definidas \textit{a priori} com base nas heurísticas de Nielsen (1994) e nos mecanismos de captura do instrumento, e as respostas foram classificadas por busca de palavras-chave e expressões associadas a cada categoria.

\section{Resultados e Discussão}

Os resultados a seguir são apresentados em três camadas complementares: (i) métricas objetivas por fluxo, (ii) evidências qualitativas, e (iii) análise exploratória de associação entre sinais objetivos e escalas UEQ-S. A discussão privilegia a interpretação por triangulação, evitando conclusões confirmatórias.

\subsection{Métricas Objetivas por Fluxo}

A Tabela 3 apresenta as métricas agregadas por fluxo para os 9 participantes. A definição e composição do \textit{FrictionScore} seguem a Seção~\ref{subsec:analise-metricas}.

\begin{table}[htbp]
  \centering
  \caption{Métricas objetivas agregadas por fluxo ($N = 9$)}
  \begin{tabular}{p{3.2cm} r r r r r}
    \toprule
    \textbf{Fluxo} & \textbf{N} & \textbf{Dur. média (s)} & \textbf{DP (s)} & \textbf{Hesit. média} & \textbf{FrictionScore} \\
    \midrule
    explorar-salvar  & 9 & 337,8 & 222,3 & 6,20 & 100,00 \\
    criar-perfil     & 9 & 253,8 & 266,9 & 2,30 &  32,02 \\
    comentar-reagir  & 9 & 162,6 & 210,0 & 3,44 &  12,29 \\
    postar-dica      & 9 & 169,5 & 135,6 & 2,22 &   2,37 \\
    \bottomrule
  \end{tabular}
\end{table}

Dois padrões aparentemente paradoxais emergem dos dados e permitem demonstrar a capacidade discriminatória do instrumento quando seus sinais são analisados em conjunto. Primeiro, o fluxo \textit{explorar-salvar} apresentou o maior FrictionScore apesar de ser baseline. A triangulação, porém, resolve a aparente contradição: a análise temática não registrou nenhuma menção de \textit{feedback} ausente, botão sem resposta ou redirecionamento falho neste fluxo, e a contagem de \textit{dead clicks} associada a ele foi a mais baixa entre todos os fluxos. Hesitações longas sem \textit{dead clicks} e sem corroboração qualitativa negativa configuram um padrão distinto, associado à leitura engajada, e não à fricção. Isso evidencia que, quando os dois sinais do instrumento (comportamento temporal e eventos discretos de clique) divergem e o qualitativo confirma a ausência de frustração, a hesitação deve ser interpretada como cognitivamente positiva. Segundo, o fluxo \textit{postar-dica}, intencionalmente degradado, apresentou o menor score de duração --- mas não o menor índice de fricção real. Os \textit{dead clicks} registrados neste fluxo, associados ao tema qualitativo ``botão sem resposta perceptível'' (33\% dos participantes), confirmam que a fricção ocorreu mesmo quando a duração bruta foi reduzida pelo comportamento de abandono precoce de parte dos participantes. A discrepância entre duração curta e frustração verbalizada sugere que, em fluxos degradados, sessões mais curtas tendem a refletir desistência, e não fluência --- distinção que o \textit{dead click} captura e a duração isolada não captura. Esses dois casos demonstram que a validade diagnóstica do instrumento emerge da leitura cruzada de seus componentes, não de qualquer métrica isolada.

Em nível de sessão, registraram-se 44 \textit{dead clicks} totais entre os 9 participantes (média: 4,9; máximo: 29, por P04), com duração média de sessão de 478,7 s (mín: 235,9 s; máx: 1011,7 s).

\begin{figure}[htbp]
  \centering
  \includegraphics[width=0.92\linewidth]{plots/flow_duration_boxplot.png}
  \caption{Distribuição da duração por fluxo (boxplot, $N = 9$).}
  \label{fig:boxplot}
\end{figure}

\subsection{Análise Temática das Respostas Qualitativas}

A Figura~\ref{fig:temas} e a Tabela 4 apresentam a frequência de cada categoria detectada na análise de conteúdo \cite{bardin1977}, conforme procedimento descrito na Seção~\ref{subsec:analise-metricas}.

\begin{figure}[htbp]
  \centering
  \includegraphics[width=0.98\linewidth]{plots/qualitative_themes_bar.png}
  \caption{Frequência das categorias qualitativas ($N = 9$).}
  \label{fig:temas}
\end{figure}

\begin{table}[htbp]
  \centering
  \caption{Temas nas respostas qualitativas ($N = 9$)}
  \begin{tabular}{p{4.5cm} r r p{4.5cm}}
    \toprule
    \textbf{Tema} & \textbf{N} & \textbf{\%} & \textbf{Relação com métrica objetiva} \\
    \midrule
    Ausência de \textit{feedback} visual & 7 & 78\% & \textit{dead clicks} e hesitações longas \\
    Botão sem resposta perceptível       & 3 & 33\% & \textit{dead clicks} $> 0$ \\
    Ausência de redirecionamento         & 3 & 33\% & hesitações longas pós-ação \\
    Lentidão / travamento percebido     & 3 & 33\% & duração elevada de fluxo \\
    Melhoria visual solicitada           & 3 & 33\% & -- \\
    Avaliação geral positiva             & 3 & 33\% & -- \\
    Confusão de navegação                & 2 & 22\% & duração elevada de fluxo \\
    Dados inconsistentes (mock)          & 2 & 22\% & -- (limitação da plataforma) \\
    \bottomrule
  \end{tabular}
\end{table}

O tema mais prevalente --- ausência de \textit{feedback} visual (78\%) --- mapeia diretamente os mecanismos de captura do instrumento. O participante P07 descreveu com precisão o fenômeno detectado automaticamente: \textit{``Ao clicar em `Salvar', fiquei esperando um retorno. Como não houve, cliquei várias vezes, entendendo que não havia funcionado anteriormente.''} De modo convergente, P08 relatou: \textit{``eu pensei que tinha travado a aplicação pois demorou e não mostrou algo diferente na tela''}. Esta triangulação confirma que os \textit{dead clicks} capturados pelo instrumento correspondem a comportamentos que os próprios usuários reconhecem e verbalizam espontaneamente.

\begin{figure}[htbp]
  \centering
  \includegraphics[width=0.92\linewidth]{plots/dead_clicks_by_participant.png}
  \caption{Distribuição de \textit{dead clicks} por participante ($N = 9$).}
  \label{fig:deadclicks}
\end{figure}

\subsection{Correlação de Spearman e Análise de Poder}

A correlação entre o \textit{FrictionScore} de sessão (Seção~\ref{subsec:analise-metricas}) e a qualidade pragmática UEQ-S resultou em $\rho = +0{,}13$ ($p \approx 0{,}73$, $N = 9$); para a escala hedônica, $\rho = -0{,}24$ ($p \approx 0{,}51$). Nenhuma correlação atingiu significância estatística. A direção negativa para a escala hedônica é teoricamente esperada: maior atrito comportamental deveria se associar a menor percepção de qualidade hedônica, uma vez que episódios de frustração pragmática contaminam a avaliação holística da experiência \cite{hassenzahl2003}.

Como o estudo é uma prova de conceito, e dado o sub-dimensionamento amostral, os coeficientes de Spearman devem ser interpretados como análise exploratória. A análise de poder \textit{a priori} (Seção~\ref{subsec:analise-metricas}) indica a necessidade de $N \geq 30$ para resultados confirmatórios.

Ainda assim, o participante P04 --- o único com UEQ-S pragmático negativo ($-0{,}25$) --- apresentou o maior FrictionScore de sessão (70,76), com 29 \textit{dead clicks} e 20 hesitações longas. Seu relato qualitativo confirmou: \textit{``Os botões não funcionam, demoram a responder, não reencaminha para o feed''}. Esse caso demonstra que, em cenários de fricção grave, o instrumento e o autorrelato convergem de forma inequívoca.

\begin{figure}[htbp]
  \centering
  \includegraphics[width=0.85\linewidth]{plots/ueqs_scatter.png}
  \caption{\textit{FrictionScore} de sessão \textit{versus} UEQ-S pragmático por participante ($N = 9$).}
  \label{fig:scatter}
\end{figure}

\subsection{Divergência Quali-Quanti e seus Significados}

Dois participantes apresentaram métricas objetivas elevadas combinadas a avaliações UEQ-S positivas: P03 (15 hesitações longas, duração de 1011 s) e P05 (44 hesitações longas). Esse padrão é consistente com dois fenômenos documentados na literatura de QoE. Primeiro, a \textit{Peak-End Rule} \cite{kahneman1993}: avaliações subjetivas de experiências são dominadas pelo pico emocional e pelo desfecho, não pela integração de todos os momentos. Quando a tarefa foi concluída com sucesso, episódios de espera intermediários tendem a ser subestimados no relato post-hoc. Segundo, o efeito de teto (\textit{ceiling effect}) do UEQ-S na versão curta: com apenas 4 itens pragmáticos e escala $-3$ a $+3$, o instrumento não discrimina diferenças dentro da faixa positiva \cite{moller2011,sauro2012}.

Esta divergência não enfraquece o estudo: pelo contrário, ela articula o argumento central em favor da medição comportamental automatizada. O instrumento captura atrito \textit{in situ} e de forma contínua, enquanto escalas subjetivas de aplicação única e poucos itens são estruturalmente limitadas para discernir fricções pontuais que o usuário supera e posteriormente racionaliza como aceitáveis.

\section{Conclusões e Trabalhos Futuros}

Encerrando a argumentação, esta seção sintetiza o que a prova de conceito permite afirmar, explicita limitações e aponta extensões diretas do instrumento e do protocolo experimental.

Este trabalho apresentou um instrumento de medição automatizada de QoE/UX baseado em Playwright, posicionado como auditor passivo de interações reais, e reportou uma prova de conceito em estudo controlado de métodos mistos. Os resultados oferecem evidências iniciais de coerência entre logs e percepção: 78\% dos participantes reportaram espontaneamente ausência de \textit{feedback} visual, em consonância com a detecção de \textit{dead clicks}. Observou-se, ainda, convergência inequívoca em um caso extremo (P04), no qual o maior FrictionScore de sessão coincidiu com UEQ-S pragmático negativo e relato explícito de botões sem resposta. Do ponto de vista metodológico, o estudo também delimita limitações: hesitações longas, quando analisadas isoladamente, permanecem ambíguas (leitura vs. desorientação) e requerem triangulação com \textit{dead clicks} e evidências qualitativas; além disso, os resultados de Spearman devem permanecer exploratórios sob $N = 9$, com indicação de $N \geq 30$ para conclusões confirmatórias. Por fim, o efeito de teto do UEQ-S na versão curta limita a discriminação dentro da faixa positiva, reforçando a complementaridade com métricas objetivas.

Como trabalhos futuros, propõe-se: (i) ampliar a amostra para $N \geq 30$ para análises confirmatórias; (ii) incorporar classificação semântica de hesitações (cognitiva vs. fricção) via análise de \textit{rawEvents}; (iii) como extensão, investigar indicadores adicionais de retrabalho técnico (por exemplo, eventos de submissão repetida e padrões de navegação de retorno) e sua eventual relação com tráfego redundante.


% ==========================================
% REFERÊNCIAS
% ==========================================
\bibliographystyle{sbc}

\bibliography{artigo}

\end{document}